\input{preamble/preamble.tex}

\geometry{margin=0.85in}
\AtBeginDocument{\small}

\newcommand{\ExamNameField}{\noindent\textbf{Name:}\ \rule{0.7\linewidth}{0.4pt}\par\medskip}

\begin{document}
	\ExamTitleBlock{10th grade}{Learning evidence 2.1 Decision analysis supplementary solutions 1}{}
	
	\ExamSection{Problems}
	\begin{ExamProblems}
		\item
		\subsection*{Problem description}
		A weekend coffee cart must choose a service plan for a street market. Plan A is a specialty menu and Plan B is a classic menu. Customer turnout can be high or low with probabilities $0.55$ and $0.45$. Expected orders are 180 for Plan A and 150 for Plan B under high turnout, and 90 for Plan A and 80 for Plan B under low turnout. The selling price per order is $8$ for Plan A and $6$ for Plan B. Variable cost per order is $3$ for Plan A and $2$ for Plan B, and fixed costs are $200$ for Plan A and $120$ for Plan B. Which decision is the risky choice under Maximax, the safe choice under Maximin, and the balanced choice under Expected Value?

		\subsection*{C2}
		\begin{center}
			\begin{tabular}{l p{0.74\linewidth}}
				\toprule
				Decision Alternatives & Plan A (specialty menu), Plan B (classic menu) \\
				States of Nature & Turnout high or low with probabilities $0.55$ and $0.45$ \\
				Events & Realized customer turnout in the market day \\
				Consequences & Profit in dollars after revenue and costs \\
				\bottomrule
			\end{tabular}
		\end{center}

		\subsection*{C3}
		Payoff = revenue $-$ cost, where revenue = (orders $\times$ price) and cost = (orders $\times$ variable cost) $+$ fixed cost.
		\begin{center}
			\textit{Orders by plan and turnout} \\
			\begin{tabular}{l c c}
				\toprule
				Plan & High turnout & Low turnout \\
				\midrule
				Plan A & 180 & 90 \\
				Plan B & 150 & 80 \\
				\bottomrule
		\end{tabular}
		\end{center}
		\begin{center}
			\begin{minipage}[t]{0.48\linewidth}
				\textit{Price parameters by plan}
				\begin{center}
					\begin{tabular}{l c}
						\toprule
						Plan & Price per order \\
						\midrule
						Plan A & 8 \\
						Plan B & 6 \\
						\bottomrule
					\end{tabular}
				\end{center}
			\end{minipage}
			\hfill
			\begin{minipage}[t]{0.48\linewidth}
				\textit{Cost parameters by plan}
				\begin{center}
					\begin{tabular}{l c c}
						\toprule
						Plan & Variable cost & Fixed cost \\
						\midrule
						Plan A & 3 & 200 \\
						Plan B & 2 & 120 \\
						\bottomrule
					\end{tabular}
				\end{center}
			\end{minipage}
		\end{center}
		\begin{center}
			\textit{Revenue by turnout} \\
			\begin{tabular}{l c c}
				\toprule
				State of nature & Plan A revenue & Plan B revenue \\
				\midrule
				High turnout &
				$\begin{array}{l}
					180\cdot 8\\
					= 1440
				\end{array}$ &
				$\begin{array}{l}
					150\cdot 6\\
					= 900
				\end{array}$ \\
				Low turnout &
				$\begin{array}{l}
					90\cdot 8\\
					= 720
				\end{array}$ &
				$\begin{array}{l}
					80\cdot 6\\
					= 480
				\end{array}$ \\
				\bottomrule
		\end{tabular}
		\end{center}
		\begin{center}
			\textit{Cost by turnout} \\
			\begin{tabular}{l c c}
				\toprule
				State of nature & Plan A cost & Plan B cost \\
				\midrule
				High turnout &
				$\begin{array}{l}
					180\cdot 3 + 200\\
					= 740
				\end{array}$ &
				$\begin{array}{l}
					150\cdot 2 + 120\\
					= 420
				\end{array}$ \\
				Low turnout &
				$\begin{array}{l}
					90\cdot 3 + 200\\
					= 470
				\end{array}$ &
				$\begin{array}{l}
					80\cdot 2 + 120\\
					= 280
				\end{array}$ \\
				\bottomrule
		\end{tabular}
		\end{center}
		Profit is computed as Revenue minus Cost for each alternative and state.
		\begin{center}
			\begin{tabular}{l c c c}
				\toprule
				State of nature & Probability & Plan A & Plan B \\
				\midrule
				High turnout & 0.55 &
				$\begin{array}{l}
					1440 - 740\\
					= 700
				\end{array}$ &
				$\begin{array}{l}
					900 - 420\\
					= 480
				\end{array}$ \\
				Low turnout & 0.45 &
				$\begin{array}{l}
					720 - 470\\
					= 250
				\end{array}$ &
				$\begin{array}{l}
					480 - 280\\
					= 200
				\end{array}$ \\
				\bottomrule
		\end{tabular}
		\end{center}

		\subsection*{C4}
		\[
		\begin{aligned}
		\max(\text{Plan A}) &= \max\{700,250\} = 700,\\
		\max(\text{Plan B}) &= \max\{480,200\} = 480.
		\end{aligned}
		\]
		\[
		\max\{\max(\text{Plan A}), \max(\text{Plan B})\}
		= \max\{700, 480\}
		= 700.
		\]
		The Maximax choice is Plan A with 700. The criterion is \emph{risky} because it focuses only on the best possible payoff and ignores how likely it is.

		\subsection*{C5}
		\[
		\begin{aligned}
		\min(\text{Plan A}) &= \min\{700,250\} = 250,\\
		\min(\text{Plan B}) &= \min\{480,200\} = 200.
		\end{aligned}
		\]
		\[
		\max\{\min(\text{Plan A}), \min(\text{Plan B})\}
		= \max\{250, 200\}
		= 250.
		\]
		The Maximin choice is Plan A with 250 because it has the least severe worst case. This is \emph{safe} because it protects against the worst outcome.

		\subsection*{C1}
		\[
		\begin{aligned}
		EV_A &= 0.55(700)+0.45(250)=385+112.5=497.5,\\
		EV_B &= 0.55(480)+0.45(200)=264+90=354.
		\end{aligned}
		\]
		The expected value criterion chooses Plan A because $497.5$ is the largest expected profit. This aligns with the Maximax and Maximin outcomes for Plan A, so it is the balanced decision using the probabilities.

		\item
		\subsection*{Problem description}
		A campus bookstore must choose a stocking plan for a semester. Plan A is premium notebooks, Plan B is standard notebooks, and Plan C is discount notebooks. Demand for notebooks can be strong, moderate, or weak with probabilities $0.35$, $0.45$, and $0.20$. Expected unit sales depend on demand and plan. Under strong demand Plan A sells 420 units, Plan B sells 520 units, and Plan C sells 600 units. Under moderate demand Plan A sells 300 units, Plan B sells 380 units, and Plan C sells 450 units. Under weak demand Plan A sells 180 units, Plan B sells 240 units, and Plan C sells 300 units. Unit prices are $22$, $18$, and $15$ for Plans A, B, and C. Variable costs per unit are $10$, $8$, and $6$, and fixed costs are $1500$, $1200$, and $900$. Which decision is the risky choice under Maximax, the safe choice under Maximin, and the balanced choice under Expected Value?

		\subsection*{C2}
		\begin{center}
			\begin{tabular}{l p{0.74\linewidth}}
				\toprule
				Decision Alternatives & Plan A (premium), Plan B (standard), Plan C (discount) \\
				States of Nature & Demand strong, moderate, weak with probabilities $0.35$, $0.45$, $0.20$ \\
				Events & Realized demand level during the semester \\
				Consequences & Profit in dollars after revenue and costs \\
				\bottomrule
			\end{tabular}
		\end{center}

		\subsection*{C3}
		Payoff = revenue $-$ cost, where revenue = (units $\times$ price) and cost = (units $\times$ variable cost) $+$ fixed cost.
		\begin{center}
			\textit{Units by plan and demand} \\
			\begin{tabular}{l c c c}
				\toprule
				Plan & Strong & Moderate & Weak \\
				\midrule
				Plan A & 420 & 300 & 180 \\
				Plan B & 520 & 380 & 240 \\
				Plan C & 600 & 450 & 300 \\
				\bottomrule
		\end{tabular}
		\end{center}
		\begin{center}
			\begin{minipage}[t]{0.48\linewidth}
				\textit{Price parameters by plan}
				\begin{center}
					\begin{tabular}{l c}
						\toprule
						Plan & Price per unit \\
						\midrule
						Plan A & 22 \\
						Plan B & 18 \\
						Plan C & 15 \\
						\bottomrule
					\end{tabular}
				\end{center}
			\end{minipage}
			\hfill
			\begin{minipage}[t]{0.48\linewidth}
				\textit{Cost parameters by plan}
				\begin{center}
					\begin{tabular}{l c c}
						\toprule
						Plan & Variable cost & Fixed cost \\
						\midrule
						Plan A & 10 & 1500 \\
						Plan B & 8 & 1200 \\
						Plan C & 6 & 900 \\
						\bottomrule
					\end{tabular}
				\end{center}
			\end{minipage}
		\end{center}
		\begin{center}
			\textit{Revenue by demand} \\
			\begin{tabular}{l c c c}
				\toprule
				State of nature & Plan A revenue & Plan B revenue & Plan C revenue \\
				\midrule
				Strong demand &
				$\begin{array}{l}
					420\cdot 22\\
					= 9240
				\end{array}$ &
				$\begin{array}{l}
					520\cdot 18\\
					= 9360
				\end{array}$ &
				$\begin{array}{l}
					600\cdot 15\\
					= 9000
				\end{array}$ \\
				Moderate demand &
				$\begin{array}{l}
					300\cdot 22\\
					= 6600
				\end{array}$ &
				$\begin{array}{l}
					380\cdot 18\\
					= 6840
				\end{array}$ &
				$\begin{array}{l}
					450\cdot 15\\
					= 6750
				\end{array}$ \\
				Weak demand &
				$\begin{array}{l}
					180\cdot 22\\
					= 3960
				\end{array}$ &
				$\begin{array}{l}
					240\cdot 18\\
					= 4320
				\end{array}$ &
				$\begin{array}{l}
					300\cdot 15\\
					= 4500
				\end{array}$ \\
				\bottomrule
		\end{tabular}
		\end{center}
		\begin{center}
			\textit{Cost by demand} \\
			\begin{tabular}{l c c c}
				\toprule
				State of nature & Plan A cost & Plan B cost & Plan C cost \\
				\midrule
				Strong demand &
				$\begin{array}{l}
					420\cdot 10 + 1500\\
					= 5700
				\end{array}$ &
				$\begin{array}{l}
					520\cdot 8 + 1200\\
					= 6560
				\end{array}$ &
				$\begin{array}{l}
					600\cdot 6 + 900\\
					= 4500
				\end{array}$ \\
				Moderate demand &
				$\begin{array}{l}
					300\cdot 10 + 1500\\
					= 4500
				\end{array}$ &
				$\begin{array}{l}
					380\cdot 8 + 1200\\
					= 4240
				\end{array}$ &
				$\begin{array}{l}
					450\cdot 6 + 900\\
					= 3600
				\end{array}$ \\
				Weak demand &
				$\begin{array}{l}
					180\cdot 10 + 1500\\
					= 3300
				\end{array}$ &
				$\begin{array}{l}
					240\cdot 8 + 1200\\
					= 3120
				\end{array}$ &
				$\begin{array}{l}
					300\cdot 6 + 900\\
					= 2700
				\end{array}$ \\
				\bottomrule
		\end{tabular}
		\end{center}
		Profit is computed as Revenue minus Cost for each alternative and state.
		\begin{center}
			\begin{tabular}{l c c c c}
				\toprule
				State of nature & Probability & Plan A & Plan B & Plan C \\
				\midrule
				Strong demand & 0.35 &
				$\begin{array}{l}
					9240 - 5700\\
					= 3540
				\end{array}$ &
				$\begin{array}{l}
					9360 - 6560\\
					= 2800
				\end{array}$ &
				$\begin{array}{l}
					9000 - 4500\\
					= 4500
				\end{array}$ \\
				Moderate demand & 0.45 &
				$\begin{array}{l}
					6600 - 4500\\
					= 2100
				\end{array}$ &
				$\begin{array}{l}
					6840 - 4240\\
					= 2600
				\end{array}$ &
				$\begin{array}{l}
					6750 - 3600\\
					= 3150
				\end{array}$ \\
				Weak demand & 0.20 &
				$\begin{array}{l}
					3960 - 3300\\
					= 660
				\end{array}$ &
				$\begin{array}{l}
					4320 - 3120\\
					= 1200
				\end{array}$ &
				$\begin{array}{l}
					4500 - 2700\\
					= 1800
				\end{array}$ \\
				\bottomrule
		\end{tabular}
		\end{center}

		\subsection*{C4}
		\[
		\begin{aligned}
		\max(\text{Plan A}) &= \max\{3540,2100,660\} = 3540,\\
		\max(\text{Plan B}) &= \max\{2800,2600,1200\} = 2800,\\
		\max(\text{Plan C}) &= \max\{4500,3150,1800\} = 4500.
		\end{aligned}
		\]
		\[
		\max\{\max(\text{Plan A}), \max(\text{Plan B}), \max(\text{Plan C})\}
		= \max\{3540, 2800, 4500\}
		= 4500.
		\]
		The Maximax choice is Plan C with 4500. The criterion is \emph{risky} because it focuses only on the best possible payoff and ignores how likely it is.

		\subsection*{C5}
		\[
		\begin{aligned}
		\min(\text{Plan A}) &= \min\{3540,2100,660\} = 660,\\
		\min(\text{Plan B}) &= \min\{2800,2600,1200\} = 1200,\\
		\min(\text{Plan C}) &= \min\{4500,3150,1800\} = 1800.
		\end{aligned}
		\]
		\[
		\max\{\min(\text{Plan A}), \min(\text{Plan B}), \min(\text{Plan C})\}
		= \max\{660, 1200, 1800\}
		= 1800.
		\]
		The Maximin choice is Plan C with 1800 because it gives the best worst case outcome. This is \emph{safe} because it prioritizes the least harmful result.

		\subsection*{C1}
		\[
		\begin{aligned}
		EV_A &= 0.35(3540)+0.45(2100)+0.20(660)=1239+945+132=2316,\\
		EV_B &= 0.35(2800)+0.45(2600)+0.20(1200)=980+1170+240=2390,\\
		EV_C &= 0.35(4500)+0.45(3150)+0.20(1800)=1575+1417.5+360=3352.5.
		\end{aligned}
		\]
		The expected value criterion chooses Plan C because $3352.5$ is the largest expected profit. This aligns with the Maximax and Maximin outcomes for Plan C, so it is the balanced decision using the probabilities.

		\item
		\subsection*{Problem description}
		A rideshare cooperative must choose a service focus for the next season. Strategy A is an urban focus and Strategy B is a suburban focus. Demand can be strong or weak with probabilities $0.65$ and $0.35$. Energy costs can be low or high with probabilities $0.55$ and $0.45$, and these costs are independent of demand. Revenue in thousands of dollars depends on demand and strategy. Under strong demand revenue is 520 for Strategy A and 560 for Strategy B. Under weak demand revenue is 260 for Strategy A and 240 for Strategy B. Operating costs in thousands of dollars depend on energy costs and strategy. Under low costs operating costs are 200 for Strategy A and 230 for Strategy B. Under high costs operating costs are 260 for Strategy A and 300 for Strategy B. Which decision is the risky choice under Maximax, the safe choice under Maximin, and the balanced choice under Expected Value?

		\subsection*{C2}
		\begin{center}
			\begin{tabular}{l p{0.74\linewidth}}
				\toprule
				Decision Alternatives & Strategy A (urban), Strategy B (suburban) \\
				States of Nature & Demand strong or weak with energy costs low or high \\
				Events & Realized demand level paired with energy cost conditions \\
				Consequences & Profit in thousands of dollars from revenue minus operating costs \\
				\bottomrule
			\end{tabular}
		\end{center}

		\subsection*{C3}
		Payoff = revenue $-$ cost, where revenue is determined by demand and cost is determined by energy prices.
		\begin{center}
			\begin{minipage}[t]{0.48\linewidth}
				\textit{Revenue parameters by strategy}
				\begin{center}
					\begin{tabular}{l c c}
						\toprule
						Strategy & Strong & Weak \\
						\midrule
						Probabilities & 0.65 & 0.35 \\
						Strategy A & 520 & 260 \\
						Strategy B & 560 & 240 \\
						\bottomrule
					\end{tabular}
				\end{center}
			\end{minipage}
			\hfill
			\begin{minipage}[t]{0.48\linewidth}
				\textit{Cost parameters by strategy}
				\begin{center}
					\begin{tabular}{l c c}
						\toprule
						Strategy & Low cost & High cost \\
						\midrule
						Probabilities & 0.55 & 0.45 \\
						Strategy A & 200 & 260 \\
						Strategy B & 230 & 300 \\
						\bottomrule
					\end{tabular}
				\end{center}
			\end{minipage}
		\end{center}
		Profit is computed as Revenue minus Cost for each alternative and state.
		\begin{center}
			\textit{Profit parameters by strategy}
			\begin{tabular}{l p{0.17\linewidth} p{0.2\linewidth} p{0.2\linewidth}}
				\toprule
				State of nature & Probability & Strategy A & Strategy B \\
				\midrule
				Strong demand, low cost &
				$\begin{array}{l}
					0.65\cdot 0.55\\
					= 0.3575
				\end{array}$ &
				$\begin{array}{l}
					520 - 200\\
					= 320
				\end{array}$ &
				$\begin{array}{l}
					560 - 230\\
					= 330
				\end{array}$ \\
				Strong demand, high cost &
				$\begin{array}{l}
					0.65\cdot 0.45\\
					= 0.2925
				\end{array}$ &
				$\begin{array}{l}
					520 - 260\\
					= 260
				\end{array}$ &
				$\begin{array}{l}
					560 - 300\\
					= 260
				\end{array}$ \\
				Weak demand, low cost &
				$\begin{array}{l}
					0.35\cdot 0.55\\
					= 0.1925
				\end{array}$ &
				$\begin{array}{l}
					260 - 200\\
					= 60
				\end{array}$ &
				$\begin{array}{l}
					240 - 230\\
					= 10
				\end{array}$ \\
				Weak demand, high cost &
				$\begin{array}{l}
					0.35\cdot 0.45\\
					= 0.1575
				\end{array}$ &
				$\begin{array}{l}
					260 - 260\\
					= 0
				\end{array}$ &
				$\begin{array}{l}
					240 - 300\\
					= -60
				\end{array}$ \\
				\bottomrule
		\end{tabular}
		\end{center}
		\begin{center}
			\textit{Payoff table} \\
			\begin{tabular}{l c c c}
				\toprule
				State of nature & Probability & Strategy A & Strategy B \\
				\midrule
				Strong demand, low cost & 0.3575 & 320 & 330 \\
				Strong demand, high cost & 0.2925 & 260 & 260 \\
				Weak demand, low cost & 0.1925 & 60 & 10 \\
				Weak demand, high cost & 0.1575 & 0 & -60 \\
				\bottomrule
		\end{tabular}
		\end{center}

		\subsection*{C4}
		\[
		\begin{aligned}
		\max(\text{Strategy A}) &= \max\{320,260,60,0\} = 320,\\
		\max(\text{Strategy B}) &= \max\{330,260,10,-60\} = 330.
		\end{aligned}
		\]
		\[
		\max\{\max(\text{Strategy A}), \max(\text{Strategy B})\}
		= \max\{320, 330\}
		= 330.
		\]
		The Maximax choice is Strategy B with 330. The criterion is \emph{risky} because it focuses only on the best possible payoff and ignores how likely it is.

		\subsection*{C5}
		\[
		\begin{aligned}
		\min(\text{Strategy A}) &= \min\{320,260,60,0\} = 0,\\
		\min(\text{Strategy B}) &= \min\{330,260,10,-60\} = -60.
		\end{aligned}
		\]
		\[
		\max\{\min(\text{Strategy A}), \min(\text{Strategy B})\}
		= \max\{0, -60\}
		= 0.
		\]
		The Maximin choice is Strategy A with 0 because it has the least severe worst case. This is \emph{safe} because it protects against the worst outcome.

		\subsection*{C1}
		Expected values use the combined state probabilities.
		\[
		\begin{aligned}
		EV_A &= 0.3575(320)+0.2925(260)+0.1925(60)+0.1575(0)\\
		&= 114.4+76.05+11.55+0=202.0,\\
		EV_B &= 0.3575(330)+0.2925(260)+0.1925(10)+0.1575(-60)\\
		&= 117.975+76.05+1.925-9.45=186.5.
		\end{aligned}
		\]
		The expected value criterion chooses Strategy A because $202.0$ is the largest expected profit. This differs from the Maximax choice and matches the Maximin choice, so the expected value emphasizes weighted outcomes.

		\item
		\subsection*{Problem description}
		A regional utility must choose a generation mix for the next year. Strategy A is solar heavy, Strategy B is wind heavy, and Strategy C is balanced. Demand can be high, medium, or low with probabilities $0.25$, $0.50$, and $0.25$. Operating cost conditions can be low or high with probabilities $0.60$ and $0.40$, and these conditions are independent of demand. Revenue in thousands of dollars depends on demand and strategy. Under high demand revenue is 680 for Strategy A, 620 for Strategy B, and 650 for Strategy C. Under medium demand revenue is 530 for Strategy A, 500 for Strategy B, and 515 for Strategy C. Under low demand revenue is 360 for Strategy A, 380 for Strategy B, and 370 for Strategy C. Operating costs in thousands of dollars depend on cost conditions and strategy. Under low costs operating costs are 300 for Strategy A, 280 for Strategy B, and 290 for Strategy C. Under high costs operating costs are 380 for Strategy A, 360 for Strategy B, and 370 for Strategy C. Which decision is the risky choice under Maximax, the safe choice under Maximin, and the balanced choice under Expected Value?

		\subsection*{C2}
		\begin{center}
			\begin{tabular}{l p{0.74\linewidth}}
				\toprule
				Decision Alternatives & Strategy A (solar), Strategy B (wind), Strategy C (balanced) \\
				States of Nature & Demand high, medium, low with cost conditions low or high \\
				Events & Realized demand level paired with operating cost conditions \\
				Consequences & Profit in thousands of dollars from revenue minus operating costs \\
				\bottomrule
			\end{tabular}
		\end{center}

		\subsection*{C3}
		Payoff = revenue $-$ cost, where revenue is determined by demand and cost is determined by operating conditions.
		\begin{center}
			\begin{minipage}[t]{0.48\linewidth}
				\textit{Revenue parameters by strategy}
				\begin{center}
					\begin{tabular}{l c c c}
						\toprule
						Strategy & High & Medium & Low \\
						\midrule
						Probabilities & 0.25 & 0.50 & 0.25 \\
						Strategy A & 680 & 530 & 360 \\
						Strategy B & 620 & 500 & 380 \\
						Strategy C & 650 & 515 & 370 \\
						\bottomrule
					\end{tabular}
				\end{center}
			\end{minipage}
			\hfill
			\begin{minipage}[t]{0.48\linewidth}
				\textit{Cost parameters by strategy}
				\begin{center}
					\begin{tabular}{l c c}
						\toprule
						Strategy & Low cost & High cost \\
						\midrule
						Probabilities & 0.60 & 0.40 \\
						Strategy A & 300 & 380 \\
						Strategy B & 280 & 360 \\
						Strategy C & 290 & 370 \\
						\bottomrule
					\end{tabular}
				\end{center}
			\end{minipage}
		\end{center}
		Profit is computed as Revenue minus Cost for each alternative and state.
		\begin{center}
			\textit{Profit parameters by strategy}
			\begin{tabular}{l p{0.17\linewidth} p{0.2\linewidth} p{0.2\linewidth} p{0.2\linewidth}}
				\toprule
				State of nature & Probability & Strategy A & Strategy B & Strategy C \\
				\midrule
				High demand, low cost &
				$\begin{array}{l}
					0.25\cdot 0.60\\
					= 0.15
				\end{array}$ &
				$\begin{array}{l}
					680 - 300\\
					= 380
				\end{array}$ &
				$\begin{array}{l}
					620 - 280\\
					= 340
				\end{array}$ &
				$\begin{array}{l}
					650 - 290\\
					= 360
				\end{array}$ \\
				High demand, high cost &
				$\begin{array}{l}
					0.25\cdot 0.40\\
					= 0.10
				\end{array}$ &
				$\begin{array}{l}
					680 - 380\\
					= 300
				\end{array}$ &
				$\begin{array}{l}
					620 - 360\\
					= 260
				\end{array}$ &
				$\begin{array}{l}
					650 - 370\\
					= 280
				\end{array}$ \\
				Medium demand, low cost &
				$\begin{array}{l}
					0.50\cdot 0.60\\
					= 0.30
				\end{array}$ &
				$\begin{array}{l}
					530 - 300\\
					= 230
				\end{array}$ &
				$\begin{array}{l}
					500 - 280\\
					= 220
				\end{array}$ &
				$\begin{array}{l}
					515 - 290\\
					= 225
				\end{array}$ \\
				Medium demand, high cost &
				$\begin{array}{l}
					0.50\cdot 0.40\\
					= 0.20
				\end{array}$ &
				$\begin{array}{l}
					530 - 380\\
					= 150
				\end{array}$ &
				$\begin{array}{l}
					500 - 360\\
					= 140
				\end{array}$ &
				$\begin{array}{l}
					515 - 370\\
					= 145
				\end{array}$ \\
				Low demand, low cost &
				$\begin{array}{l}
					0.25\cdot 0.60\\
					= 0.15
				\end{array}$ &
				$\begin{array}{l}
					360 - 300\\
					= 60
				\end{array}$ &
				$\begin{array}{l}
					380 - 280\\
					= 100
				\end{array}$ &
				$\begin{array}{l}
					370 - 290\\
					= 80
				\end{array}$ \\
				Low demand, high cost &
				$\begin{array}{l}
					0.25\cdot 0.40\\
					= 0.10
				\end{array}$ &
				$\begin{array}{l}
					360 - 380\\
					= -20
				\end{array}$ &
				$\begin{array}{l}
					380 - 360\\
					= 20
				\end{array}$ &
				$\begin{array}{l}
					370 - 370\\
					= 0
				\end{array}$ \\
				\bottomrule
		\end{tabular}
		\end{center}
		\begin{center}
			\textit{Payoff table} \\
			\begin{tabular}{l c c c c}
				\toprule
				State of nature & Probability & Strategy A & Strategy B & Strategy C \\
				\midrule
				High demand, low cost & 0.15 & 380 & 340 & 360 \\
				High demand, high cost & 0.10 & 300 & 260 & 280 \\
				Medium demand, low cost & 0.30 & 230 & 220 & 225 \\
				Medium demand, high cost & 0.20 & 150 & 140 & 145 \\
				Low demand, low cost & 0.15 & 60 & 100 & 80 \\
				Low demand, high cost & 0.10 & -20 & 20 & 0 \\
				\bottomrule
		\end{tabular}
		\end{center}

		\subsection*{C4}
		\[
		\begin{aligned}
		\max(\text{Strategy A}) &= \max\{380,300,230,150,60,-20\} = 380,\\
		\max(\text{Strategy B}) &= \max\{340,260,220,140,100,20\} = 340,\\
		\max(\text{Strategy C}) &= \max\{360,280,225,145,80,0\} = 360.
		\end{aligned}
		\]
		\[
		\max\{\max(\text{Strategy A}), \max(\text{Strategy B}), \max(\text{Strategy C})\}
		= \max\{380, 340, 360\}
		= 380.
		\]
		The Maximax choice is Strategy A with 380. The criterion is \emph{risky} because it focuses only on the best possible payoff and ignores how likely it is.

		\subsection*{C5}
		\[
		\begin{aligned}
		\min(\text{Strategy A}) &= \min\{380,300,230,150,60,-20\} = -20,\\
		\min(\text{Strategy B}) &= \min\{340,260,220,140,100,20\} = 20,\\
		\min(\text{Strategy C}) &= \min\{360,280,225,145,80,0\} = 0.
		\end{aligned}
		\]
		\[
		\max\{\min(\text{Strategy A}), \min(\text{Strategy B}), \min(\text{Strategy C})\}
		= \max\{-20, 20, 0\}
		= 20.
		\]
		The Maximin choice is Strategy B with 20 because it has the least severe worst case. This is \emph{safe} because it protects against the worst outcome.

		\subsection*{C1}
		Expected values use the combined state probabilities.
		\[
		\begin{aligned}
		EV_A &= 0.15(380)+0.10(300)+0.30(230)+0.20(150)+0.15(60)+0.10(-20)\\
		&= 57+30+69+30+9-2=193,\\
		EV_B &= 0.15(340)+0.10(260)+0.30(220)+0.20(140)+0.15(100)+0.10(20)\\
		&= 51+26+66+28+15+2=188,\\
		EV_C &= 0.15(360)+0.10(280)+0.30(225)+0.20(145)+0.15(80)+0.10(0)\\
		&= 54+28+67.5+29+12+0=190.5.
		\end{aligned}
		\]
		The expected value criterion chooses Strategy A because $193$ is the largest expected profit. This differs from the Maximin choice and matches the Maximax choice, so the expected value gives a balanced decision using the probabilities.

		\item
		\subsection*{Problem description}
		A community fitness hub must choose a membership model for the next year. Model A is a weekday focus, Model B is a family pass, and Model C is an evening pass. Demand can be very high, high, medium, or low with probabilities $0.20$, $0.30$, $0.30$, and $0.20$. Staffing costs can be favorable or unfavorable with probabilities $0.55$ and $0.45$, and these costs are independent of demand. Expected membership counts depend on demand and model. Under very high demand Model A has 820 members, Model B has 760 members, and Model C has 900 members. Under high demand Model A has 700 members, Model B has 650 members, and Model C has 760 members. Under medium demand Model A has 520 members, Model B has 480 members, and Model C has 600 members. Under low demand Model A has 360 members, Model B has 320 members, and Model C has 420 members. Membership fees are $38$ for Model A, $42$ for Model B, and $34$ for Model C per member per year. Variable staffing cost per member is $20$ when costs are favorable and $28$ when costs are unfavorable. Fixed annual costs are $8200$ for Model A, $9600$ for Model B, and $7000$ for Model C. Which decision is the risky choice under Maximax, the safe choice under Maximin, and the balanced choice under Expected Value?

		\subsection*{C2}
		\begin{center}
			\begin{tabular}{l p{0.74\linewidth}}
				\toprule
				Decision Alternatives & Model A (weekday), Model B (family), Model C (evening) \\
				States of Nature & Demand very high, high, medium, low with staffing costs favorable or unfavorable \\
				Events & Realized demand level paired with staffing cost conditions \\
				Consequences & Annual profit in dollars from membership fees minus staffing and fixed costs \\
				\bottomrule
			\end{tabular}
		\end{center}

		\subsection*{C3}
		Payoff = revenue $-$ cost, where revenue = (members $\times$ fee) and cost = (members $\times$ variable cost) $+$ fixed cost.
		\begin{center}
			\textit{Members per model by demand} \\
			\begin{tabular}{l c c c c}
				\toprule
				Model & Very high & High & Medium & Low \\
				\midrule
				Model A & 820 & 700 & 520 & 360 \\
				Model B & 760 & 650 & 480 & 320 \\
				Model C & 900 & 760 & 600 & 420 \\
				\bottomrule
		\end{tabular}
		\end{center}
		\begin{center}
			\begin{minipage}[t]{0.48\linewidth}
				\textit{Cost parameters by model}
				\begin{center}
					\begin{tabular}{l c c c}
						\toprule
						Model & Variable favorable & Variable unfavorable & Fixed cost \\
						\midrule
						Model A & 20 & 28 & 8200 \\
						Model B & 20 & 28 & 9600 \\
						Model C & 20 & 28 & 7000 \\
						\bottomrule
					\end{tabular}
				\end{center}
			\end{minipage}
			\hfill
			\begin{minipage}[t]{0.48\linewidth}
				\textit{Revenue parameters by member}
				\begin{center}
					\begin{tabular}{l c}
						\toprule
						Model & Fee per member \\
						\midrule
						Model A & 38 \\
						Model B & 42 \\
						Model C & 34 \\
						\bottomrule
					\end{tabular}
				\end{center}
			\end{minipage}
		\end{center}
		\begin{center}
			\textit{Revenue by demand} \\
			\begin{tabular}{l c c c}
				\toprule
				State of nature & Model A revenue & Model B revenue & Model C revenue \\
				\midrule
				Very high demand &
				$\begin{array}{l}
					820\cdot 38\\
					= 31160
				\end{array}$ &
				$\begin{array}{l}
					760\cdot 42\\
					= 31920
				\end{array}$ &
				$\begin{array}{l}
					900\cdot 34\\
					= 30600
				\end{array}$ \\
				High demand &
				$\begin{array}{l}
					700\cdot 38\\
					= 26600
				\end{array}$ &
				$\begin{array}{l}
					650\cdot 42\\
					= 27300
				\end{array}$ &
				$\begin{array}{l}
					760\cdot 34\\
					= 25840
				\end{array}$ \\
				Medium demand &
				$\begin{array}{l}
					520\cdot 38\\
					= 19760
				\end{array}$ &
				$\begin{array}{l}
					480\cdot 42\\
					= 20160
				\end{array}$ &
				$\begin{array}{l}
					600\cdot 34\\
					= 20400
				\end{array}$ \\
				Low demand &
				$\begin{array}{l}
					360\cdot 38\\
					= 13680
				\end{array}$ &
				$\begin{array}{l}
					320\cdot 42\\
					= 13440
				\end{array}$ &
				$\begin{array}{l}
					420\cdot 34\\
					= 14280
				\end{array}$ \\
				\bottomrule
		\end{tabular}
		\end{center}
		\begin{center}
			\textit{Cost by demand and staffing condition} \\
			\begin{tabular}{l c c c}
				\toprule
				State of nature & Model A cost & Model B cost & Model C cost \\
				\midrule
				Very high demand, favorable cost &
				$\begin{array}{l}
					820\cdot 20 + 8200\\
					= 24600
				\end{array}$ &
				$\begin{array}{l}
					760\cdot 20 + 9600\\
					= 24800
				\end{array}$ &
				$\begin{array}{l}
					900\cdot 20 + 7000\\
					= 25000
				\end{array}$ \\
				Very high demand, unfavorable cost &
				$\begin{array}{l}
					820\cdot 28 + 8200\\
					= 31160
				\end{array}$ &
				$\begin{array}{l}
					760\cdot 28 + 9600\\
					= 30880
				\end{array}$ &
				$\begin{array}{l}
					900\cdot 28 + 7000\\
					= 32200
				\end{array}$ \\
				High demand, favorable cost &
				$\begin{array}{l}
					700\cdot 20 + 8200\\
					= 22200
				\end{array}$ &
				$\begin{array}{l}
					650\cdot 20 + 9600\\
					= 22600
				\end{array}$ &
				$\begin{array}{l}
					760\cdot 20 + 7000\\
					= 22200
				\end{array}$ \\
				High demand, unfavorable cost &
				$\begin{array}{l}
					700\cdot 28 + 8200\\
					= 27800
				\end{array}$ &
				$\begin{array}{l}
					650\cdot 28 + 9600\\
					= 27800
				\end{array}$ &
				$\begin{array}{l}
					760\cdot 28 + 7000\\
					= 28280
				\end{array}$ \\
				Medium demand, favorable cost &
				$\begin{array}{l}
					520\cdot 20 + 8200\\
					= 18600
				\end{array}$ &
				$\begin{array}{l}
					480\cdot 20 + 9600\\
					= 19200
				\end{array}$ &
				$\begin{array}{l}
					600\cdot 20 + 7000\\
					= 19000
				\end{array}$ \\
				Medium demand, unfavorable cost &
				$\begin{array}{l}
					520\cdot 28 + 8200\\
					= 22760
				\end{array}$ &
				$\begin{array}{l}
					480\cdot 28 + 9600\\
					= 23040
				\end{array}$ &
				$\begin{array}{l}
					600\cdot 28 + 7000\\
					= 23800
				\end{array}$ \\
				Low demand, favorable cost &
				$\begin{array}{l}
					360\cdot 20 + 8200\\
					= 15400
				\end{array}$ &
				$\begin{array}{l}
					320\cdot 20 + 9600\\
					= 16000
				\end{array}$ &
				$\begin{array}{l}
					420\cdot 20 + 7000\\
					= 15400
				\end{array}$ \\
				Low demand, unfavorable cost &
				$\begin{array}{l}
					360\cdot 28 + 8200\\
					= 18280
				\end{array}$ &
				$\begin{array}{l}
					320\cdot 28 + 9600\\
					= 18560
				\end{array}$ &
				$\begin{array}{l}
					420\cdot 28 + 7000\\
					= 18760
				\end{array}$ \\
				\bottomrule
		\end{tabular}
		\end{center}
		Profit is computed as Revenue minus Cost for each alternative and state.
		\begin{center}
			\textit{Profit parameters by model} \\
			\begin{tabular}{l p{0.19\linewidth} p{0.19\linewidth} p{0.19\linewidth} p{0.19\linewidth}}
				\toprule
				State of nature & Model A & Model B & Model C & Probabilities \\
				\midrule
				Very high demand, favorable cost &
				$\begin{array}{l}
					31160 - 24600\\
					= 6560
				\end{array}$ &
				$\begin{array}{l}
					31920 - 24800\\
					= 7120
				\end{array}$ &
				$\begin{array}{l}
					30600 - 25000\\
					= 5600
				\end{array}$ &
				$\begin{array}{l}
					0.20\cdot 0.55\\
					= 0.11
				\end{array}$ \\
				Very high demand, unfavorable cost &
				$\begin{array}{l}
					31160 - 31160\\
					= 0
				\end{array}$ &
				$\begin{array}{l}
					31920 - 30880\\
					= 1040
				\end{array}$ &
				$\begin{array}{l}
					30600 - 32200\\
					= -1600
				\end{array}$ &
				$\begin{array}{l}
					0.20\cdot 0.45\\
					= 0.09
				\end{array}$ \\
				High demand, favorable cost &
				$\begin{array}{l}
					26600 - 22200\\
					= 4400
				\end{array}$ &
				$\begin{array}{l}
					27300 - 22600\\
					= 4700
				\end{array}$ &
				$\begin{array}{l}
					25840 - 22200\\
					= 3640
				\end{array}$ &
				$\begin{array}{l}
					0.30\cdot 0.55\\
					= 0.165
				\end{array}$ \\
				High demand, unfavorable cost &
				$\begin{array}{l}
					26600 - 27800\\
					= -1200
				\end{array}$ &
				$\begin{array}{l}
					27300 - 27800\\
					= -500
				\end{array}$ &
				$\begin{array}{l}
					25840 - 28280\\
					= -2440
				\end{array}$ &
				$\begin{array}{l}
					0.30\cdot 0.45\\
					= 0.135
				\end{array}$ \\
				Medium demand, favorable cost &
				$\begin{array}{l}
					19760 - 18600\\
					= 1160
				\end{array}$ &
				$\begin{array}{l}
					20160 - 19200\\
					= 960
				\end{array}$ &
				$\begin{array}{l}
					20400 - 19000\\
					= 1400
				\end{array}$ &
				$\begin{array}{l}
					0.30\cdot 0.55\\
					= 0.165
				\end{array}$ \\
				Medium demand, unfavorable cost &
				$\begin{array}{l}
					19760 - 22760\\
					= -3000
				\end{array}$ &
				$\begin{array}{l}
					20160 - 23040\\
					= -2880
				\end{array}$ &
				$\begin{array}{l}
					20400 - 23800\\
					= -3400
				\end{array}$ &
				$\begin{array}{l}
					0.30\cdot 0.45\\
					= 0.135
				\end{array}$ \\
				Low demand, favorable cost &
				$\begin{array}{l}
					13680 - 15400\\
					= -1720
				\end{array}$ &
				$\begin{array}{l}
					13440 - 16000\\
					= -2560
				\end{array}$ &
				$\begin{array}{l}
					14280 - 15400\\
					= -1120
				\end{array}$ &
				$\begin{array}{l}
					0.20\cdot 0.55\\
					= 0.11
				\end{array}$ \\
				Low demand, unfavorable cost &
				$\begin{array}{l}
					13680 - 18280\\
					= -4600
				\end{array}$ &
				$\begin{array}{l}
					13440 - 18560\\
					= -5120
				\end{array}$ &
				$\begin{array}{l}
					14280 - 18760\\
					= -4480
				\end{array}$ &
				$\begin{array}{l}
					0.20\cdot 0.45\\
					= 0.09
				\end{array}$ \\
				\bottomrule
		\end{tabular}
		\end{center}
		Combined probabilities are $0.20\cdot 0.55=0.11$, $0.20\cdot 0.45=0.09$, $0.30\cdot 0.55=0.165$, $0.30\cdot 0.45=0.135$, $0.30\cdot 0.55=0.165$, $0.30\cdot 0.45=0.135$, $0.20\cdot 0.55=0.11$, $0.20\cdot 0.45=0.09$.
		\begin{center}
			\begin{tabular}{l c c c c}
				\toprule
				State of nature & Probability & Model A & Model B & Model C \\
				\midrule
				Very high demand, favorable cost & 0.11 & 6560 & 7120 & 5600 \\
				Very high demand, unfavorable cost & 0.09 & 0 & 1040 & -1600 \\
				High demand, favorable cost & 0.165 & 4400 & 4700 & 3640 \\
				High demand, unfavorable cost & 0.135 & -1200 & -500 & -2440 \\
				Medium demand, favorable cost & 0.165 & 1160 & 960 & 1400 \\
				Medium demand, unfavorable cost & 0.135 & -3000 & -2880 & -3400 \\
				Low demand, favorable cost & 0.11 & -1720 & -2560 & -1120 \\
				Low demand, unfavorable cost & 0.09 & -4600 & -5120 & -4480 \\
				\bottomrule
		\end{tabular}
		\end{center}

		\subsection*{C4}
		\[
		\begin{aligned}
		\max(\text{Model A}) &= \max\{6560,0,4400,-1200,1160,-3000,-1720,-4600\} = 6560,\\
		\max(\text{Model B}) &= \max\{7120,1040,4700,-500,960,-2880,-2560,-5120\} = 7120,\\
		\max(\text{Model C}) &= \max\{5600,-1600,3640,-2440,1400,-3400,-1120,-4480\} = 5600.
		\end{aligned}
		\]
		\[
		\max\{\max(\text{Model A}), \max(\text{Model B}), \max(\text{Model C})\}
		= \max\{6560, 7120, 5600\}
		= 7120.
		\]
		The Maximax choice is Model B with 7120. The criterion is \emph{risky} because it focuses only on the best possible payoff and ignores how likely it is.

		\subsection*{C5}
		\[
		\begin{aligned}
		\min(\text{Model A}) &= \min\{6560,0,4400,-1200,1160,-3000,-1720,-4600\} = -4600,\\
		\min(\text{Model B}) &= \min\{7120,1040,4700,-500,960,-2880,-2560,-5120\} = -5120,\\
		\min(\text{Model C}) &= \min\{5600,-1600,3640,-2440,1400,-3400,-1120,-4480\} = -4480.
		\end{aligned}
		\]
		\[
		\max\{\min(\text{Model A}), \min(\text{Model B}), \min(\text{Model C})\}
		= \max\{-4600, -5120, -4480\}
		= -4480.
		\]
		The Maximin choice is Model C with $-4480$ because it has the least severe worst case loss. This is \emph{safe} because it protects against the worst outcome.

		\subsection*{C1}
		Expected values use the combined state probabilities.
		\[
		\begin{aligned}
		EV_A &= 0.11(6560)+0.09(0)+0.165(4400)+0.135(-1200)+0.165(1160)+0.135(-3000)+0.11(-1720)+0.09(-4600)\\
		&= 721.6+0+726-162+191.4-405-189.2-414=468.8,\\
		EV_B &= 0.11(7120)+0.09(1040)+0.165(4700)+0.135(-500)+0.165(960)+0.135(-2880)+0.11(-2560)+0.09(-5120)\\
		&= 783.2+93.6+775.5-67.5+158.4-388.8-281.6-460.8=612.0,\\
		EV_C &= 0.11(5600)+0.09(-1600)+0.165(3640)+0.135(-2440)+0.165(1400)+0.135(-3400)+0.11(-1120)+0.09(-4480)\\
		&= 616-144+600.6-329.4+231-459-123.2-403.2=-11.2.
		\end{aligned}
		\]
		The expected value criterion chooses Model B because $612.0$ is the largest expected profit. This differs from the Maximin choice and matches the Maximax choice, so the expected value gives a balanced decision using the probabilities.
	\end{ExamProblems}
\end{document}
